\documentclass[11pt]{beamer}
\usepackage[utf8]{inputenc}
\usepackage{amsmath, amssymb, graphicx}
\usepackage{bm} % For bold math symbols (vectors)

% Presentation Theme
\usetheme{Madrid}
\usecolortheme{default}

\title[Lambert's Problem]{Astrodynamics 2026-1}
\subtitle{Lambert's Boundary Value Problem \& Launch Energy ($C_3$)}
\author{Prof. Juan}
\institute{Universidad de Antioquia - Aerospace Engineering}
\date{Spring 2026}

\begin{document}

% --- Title Slide ---
\begin{frame}
    \titlepage
\end{frame}

% ==========================================
% SECTION 1: THE BOUNDARY VALUE PROBLEM
% ==========================================
\section{The Boundary Value Problem}

% --- Slide 1 ---
\begin{frame}{Hohmann vs. Lambert}
    \textbf{The Hohmann Transfer (The Forward Problem):}
    \begin{itemize}
        \item We constrain the \textit{geometry} (an ellipse tangent to two circular orbits).
        \item The math dictates the required \textit{Time of Flight} ($\Delta t$).
    \end{itemize}
    
    \vspace{0.3cm}
    \textbf{Lambert's Problem (The Inverse Problem):}
    \begin{itemize}
        \item We dictate the exact \textit{Time of Flight} ($\Delta t$).
        \item The math dictates the required \textit{geometry} of the transfer ellipse.
        \item This is essential for real planetary ephemerides, where planets are never perfectly aligned for a Hohmann transfer.
    \end{itemize}
\end{frame}

% --- Slide 2 ---
\begin{frame}{The Knowns and Unknowns}
    To solve Lambert's problem, we set up a Boundary Value Problem (BVP).
    
    \vspace{0.3cm}
    \textbf{The Inputs (Boundary Conditions):}
    \begin{enumerate}
        \item $\bm{r}_1$: The 3D position vector of the departure planet at $t_{dep}$.
        \item $\bm{r}_2$: The 3D position vector of the arrival planet at $t_{arr}$.
        \item $\Delta t$: The Time of Flight ($t_{arr} - t_{dep}$).
    \end{enumerate}
    
    \vspace{0.3cm}
    \textbf{The Outputs (What we must solve for):}
    \begin{itemize}
        \item $\bm{v}_1$: The required heliocentric velocity vector at departure.
        \item $\bm{v}_2$: The resulting heliocentric velocity vector at arrival.
    \end{itemize}
\end{frame}

% ==========================================
% SECTION 2: TRANSFER GEOMETRY
% ==========================================
\section{Transfer Geometry}

% --- Slide 3 ---
\begin{frame}{Defining the Transfer Plane}
    The position vectors $\bm{r}_1$ and $\bm{r}_2$ extend outward from the Sun. Together, they define a specific 2D plane in space (the transfer plane).
    
    \vspace{0.3cm}
    \textbf{Key Geometric Parameters:}
    \begin{itemize}
        \item \textbf{Transfer Angle ($\Delta \nu$):} The angle between $\bm{r}_1$ and $\bm{r}_2$.
        \item \textbf{The Chord ($c$):} The straight-line distance connecting the two planet positions.
    \end{itemize}
    \begin{equation}
        c = \sqrt{r_1^2 + r_2^2 - 2 r_1 r_2 \cos(\Delta \nu)}
    \end{equation}
\end{frame}

% --- Slide 4 ---
\begin{frame}{The 180-Degree Singularity (The "Ridge")}
    What happens if the planets are exactly $180^\circ$ apart ($\Delta \nu = \pi$)?
    
    \vspace{0.3cm}
    \begin{itemize}
        \item The vectors $\bm{r}_1$ and $\bm{r}_2$ form a straight line, not a unique plane.
        \item There are \textit{infinite} possible planes that can connect them.
        \item If the planets have different orbital inclinations, a perfectly $180^\circ$ transfer requires an immense plane-change maneuver at departure.
        \item This mathematical singularity appears as the high-energy \textbf{"Ridge"} dividing the Type I and Type II lobes on a Porkchop Plot.
    \end{itemize}
\end{frame}

% ==========================================
% SECTION 3: LAMBERT'S THEOREM
% ==========================================
\section{Lambert's Theorem}

% --- Slide 5 ---
\begin{frame}{Lambert's Theorem}
    The numerical solvers we use (like the Izzo solver in Python) rely on a fundamental principle of orbital mechanics:
    
    \vspace{0.3cm}
    \begin{block}{Lambert's Theorem}
        The transfer time ($\Delta t$) between any two points on a conic trajectory depends \textbf{ONLY} on the sum of the radial distances ($r_1 + r_2$), the chord length ($c$), and the semi-major axis ($a$) of the transfer conic.
    \end{block}
    
    \vspace{0.3cm}
    Because $r_1$, $r_2$, and $c$ are fixed by the planetary ephemerides, the only free variable is the semi-major axis ($a$).
\end{frame}

% --- Slide 6 ---
\begin{frame}{The Iterative Solution}
    Because the relationship between time and the semi-major axis is transcendental (like Kepler's Equation), we cannot solve it algebraically.
    
    \vspace{0.3cm}
    \textbf{The Algorithmic Process:}
    \begin{enumerate}
        \item The computer guesses a value for the semi-major axis ($a$).
        \item It calculates the theoretical Time of Flight for that specific ellipse.
        \item It compares the calculated time to our desired $\Delta t$.
        \item It iterates using numerical methods (e.g., Newton-Raphson) until the times match.
        \item Once converged, it calculates the departure velocity vector ($\bm{v}_1$).
    \end{enumerate}
\end{frame}

% ==========================================
% SECTION 4: DEPARTURE ENERGY
% ==========================================
\section{Departure Energy ($C_3$)}

% --- Slide 7 ---
\begin{frame}{From Heliocentric to Geocentric}
    The solver outputs $\bm{v}_1$, the required \textit{heliocentric} velocity. We must translate this into launch vehicle requirements.
    
    \vspace{0.3cm}
    \textbf{1. Hyperbolic Excess Velocity ($\bm{v}_\infty$):} \\
    The spacecraft is already moving with the Earth ($\bm{v}_{Earth}$). The launch vehicle only needs to provide the vector difference.
    \begin{equation}
        \bm{v}_\infty = \bm{v}_1 - \bm{v}_{Earth}
    \end{equation}
    
    \vspace{0.2cm}
    \textit{Note: This $\bm{v}_\infty$ is the exact same asymptote vector we targeted in our Patched Conics GMAT lab!}
\end{frame}

% --- Slide 8 ---
\begin{frame}{Calculating Characteristic Energy ($C_3$)}
    \textbf{2. Characteristic Energy ($C_3$):} \\
    Finally, we calculate the scalar energy state required for the launch vehicle. $C_3$ is the square of the magnitude of the excess velocity vector.
    
    \begin{equation}
        C_3 = \bm{v}_\infty \cdot \bm{v}_\infty = |\bm{v}_\infty|^2
    \end{equation}
    
    \vspace{0.3cm}
    \textbf{Conclusion:} 
    This final $C_3$ value is exactly what gets plotted on the Z-axis of our Porkchop plots. By looping this entire process over thousands of dates, we find the optimal launch windows!
\end{frame}

\end{document}